% GENERATED FILE. Edit canonical lessons, then run npm run generate:books.
% canonical-source-hash: fnv1a64:c6994d2c74285e3c
% canonical-lessons: LA-C270-circus, LA-C270-basilica, LA-C270-curia, LA-C270-ancora, LA-C270-remus

\chapter{Racecourse, Hall of justice, Senate house, Anchor, Oar}
\label{ch:la-racecourse-hall-of-justice-senate-house-anchor-oar}
\hlchaptermodality{270}

\begin{chapteropening}
\textbf{By the end of this chapter:} \emph{I can say a racecourse, a hall of justice, the senate house, an anchor and an oar.}

Complete the last lesson of chapter 270: I can say a racecourse, a hall of justice, the senate house, an anchor and an oar.
\end{chapteropening}

\section[circus, circī]{circus, circī — a racecourse}

\label{lesson:LA-C270-circus}



\begin{quote}
Before the new one: say the Latin for greedy, then the Latin for thrifty.
\end{quote}

\subsection*{You'll want to know: circus, circī}
\textbf{circus, circī} — \textquotedblleft{}a racecourse\textquotedblright{}.

\begin{cousinweb}
\textbf{circus} is usually taken from Greek \emph{kirkos}, \textquotedblleft{}a ring, a circle\textquotedblright{} (Lewis \& Short), a form beside \emph{krikos}, \textquotedblleft{}ring\textquotedblright{}. English \emph{circle} comes from its diminutive \emph{circulus}, \textquotedblleft{}a small ring\textquotedblright{}.
\end{cousinweb}

\begin{grammarlens}[title={What you've built}]
One more word for this chapter.
\end{grammarlens}

\subsection*{Your turn}
\begin{itemize}
\raggedright
  \item \emph{Say it:} \emph{circus, circī}
  \item \emph{Say it:} \emph{circus, circī}, once more
  \item \emph{Say it:} say \emph{circus, circī}
  \item \emph{Recall:} say the Latin for greedy, then the Latin for thrifty, then say \emph{circus, circī} again
\end{itemize}

\begin{tcolorbox}[breakable,colback=teal!4,colframe=teal!35!black,title={Before you move on}]
What does circus, circī mean? (\textquotedblleft{}A racecourse\textquotedblright{}.) Say it once more.
\end{tcolorbox}
\section[basilica, basilicae]{basilica, basilicae — a hall of justice}

\label{lesson:LA-C270-basilica}



\begin{quote}
Before the new one: say the Latin for thrifty, then the Latin for a racecourse.
\end{quote}

\subsection*{You'll want to know: basilica, basilicae}
\textbf{basilica, basilicae} — \textquotedblleft{}a hall of justice\textquotedblright{}.

\begin{cousinweb}
\textbf{basilica} was borrowed from Greek \emph{basilikē}, \textquotedblleft{}royal\textquotedblright{}, short for a royal hall (EtymDB 2.1), from \emph{basileus}, \textquotedblleft{}king\textquotedblright{}.
\end{cousinweb}

\begin{grammarlens}[title={What you've built}]
One more word for this chapter.
\end{grammarlens}

\subsection*{Your turn}
\begin{itemize}
\raggedright
  \item \emph{Say it:} \emph{basilica, basilicae}
  \item \emph{Say it:} \emph{basilica, basilicae}, once more
  \item \emph{Say it:} say \emph{basilica, basilicae}
  \item \emph{Recall:} say the Latin for thrifty, then the Latin for a racecourse, then say \emph{basilica, basilicae} again
\end{itemize}

\begin{tcolorbox}[breakable,colback=teal!4,colframe=teal!35!black,title={Before you move on}]
What does basilica, basilicae mean? (\textquotedblleft{}A hall of justice\textquotedblright{}.) Say it once more.
\end{tcolorbox}
\section[cūria, cūriae]{cūria, cūriae — the senate house}

\label{lesson:LA-C270-curia}



\begin{quote}
Before the new one: say the Latin for a racecourse, then the Latin for a hall of justice.
\end{quote}

\subsection*{You'll want to know: cūria, cūriae}
\textbf{cūria, cūriae} — \textquotedblleft{}the senate house\textquotedblright{}.

\begin{cousinweb}
\textbf{cūria} is often explained as \emph{*co-viria}, \textquotedblleft{}a gathering of men\textquotedblright{} (EtymDB 2.1), from \emph{co-}, \textquotedblleft{}together\textquotedblright{}, and \emph{vir}, \textquotedblleft{}a man\textquotedblright{}, but this is uncertain.
\end{cousinweb}

\begin{grammarlens}[title={What you've built}]
One more word for this chapter.
\end{grammarlens}

\subsection*{Your turn}
\begin{itemize}
\raggedright
  \item \emph{Say it:} \emph{cūria, cūriae}
  \item \emph{Say it:} \emph{cūria, cūriae}, once more
  \item \emph{Say it:} say \emph{cūria, cūriae}
  \item \emph{Recall:} say the Latin for a racecourse, then the Latin for a hall of justice, then say \emph{cūria, cūriae} again
\end{itemize}

\begin{tcolorbox}[breakable,colback=teal!4,colframe=teal!35!black,title={Before you move on}]
What does cūria, cūriae mean? (\textquotedblleft{}The senate house\textquotedblright{}.) Say it once more.
\end{tcolorbox}
\section[ancora, ancorae]{ancora, ancorae — an anchor}

\label{lesson:LA-C270-ancora}



\begin{quote}
Before the new one: say the Latin for a hall of justice, then the Latin for the senate house.
\end{quote}

\subsection*{You'll want to know: ancora, ancorae}
\textbf{ancora, ancorae} — \textquotedblleft{}an anchor\textquotedblright{}.

\begin{cousinweb}
\textbf{ancora} was borrowed from Greek \emph{ankyra}, \textquotedblleft{}anchor\textquotedblright{} (EtymDB 2.1), from a root for bending: a hooked thing.
\end{cousinweb}

\begin{grammarlens}[title={What you've built}]
One more word for this chapter.
\end{grammarlens}

\subsection*{Your turn}
\begin{itemize}
\raggedright
  \item \emph{Say it:} \emph{ancora, ancorae}
  \item \emph{Say it:} \emph{ancora, ancorae}, once more
  \item \emph{Say it:} say \emph{ancora, ancorae}
  \item \emph{Recall:} say the Latin for a hall of justice, then the Latin for the senate house, then say \emph{ancora, ancorae} again
\end{itemize}

\begin{tcolorbox}[breakable,colback=teal!4,colframe=teal!35!black,title={Before you move on}]
What does ancora, ancorae mean? (\textquotedblleft{}An anchor\textquotedblright{}.) Say it once more.
\end{tcolorbox}
\section[rēmus, rēmī]{rēmus, rēmī — an oar}

\label{lesson:LA-C270-remus}



\begin{quote}
Before the new one: say the Latin for the senate house, then the Latin for an anchor.
\end{quote}

\subsection*{You'll want to know: rēmus, rēmī}
\textbf{rēmus, rēmī} — \textquotedblleft{}an oar\textquotedblright{}.

\begin{cousinweb}
\textbf{rēmus} is inherited, usually traced to Proto-Indo-European \emph{*h1reh1-}, \textquotedblleft{}to row\textquotedblright{}. Greek \emph{eretmos}, \textquotedblleft{}oar\textquotedblright{}, is a cousin (Lewis \& Short), and so are English \emph{row} and \emph{rudder}.
\end{cousinweb}

\begin{grammarlens}[title={What you've built}]
One more word for this chapter.
\end{grammarlens}

\subsection*{Your turn}
\begin{itemize}
\raggedright
  \item \emph{Say it:} \emph{rēmus, rēmī}
  \item \emph{Say it:} \emph{rēmus, rēmī}, once more
  \item \emph{Say it:} say \emph{rēmus, rēmī}
  \item \emph{Recall:} say the Latin for the senate house, then the Latin for an anchor, then say \emph{rēmus, rēmī} again
\end{itemize}

\begin{tcolorbox}[breakable,colback=teal!4,colframe=teal!35!black,title={Before you move on}]
What does rēmus, rēmī mean? (\textquotedblleft{}An oar\textquotedblright{}.) Say it once more.
\end{tcolorbox}
